\documentclass[10pt,a4paper]{amsart}
\usepackage[margin=19mm]{geometry}
\usepackage{amsmath,amssymb,mathtools}
\usepackage[scr=rsfs,cal=euler]{mathalfa}
\usepackage{xfrac}
\usepackage[unicode,hidelinks,bookmarks=false]{hyperref}
\newcommand{\ie}{i.e.\ }
\newcommand{\cf}{cf.\ }
\newcommand{\resp}{respectively\ }
\newcommand{\Subsection}{\subsection}
\providecommand{\nobreakdash}{\nobreak\hbox{-}}
\setlength{\emergencystretch}{2em}
\multlinegap0pt
\usepackage{amssymb}
\usepackage{xcolor}
\usepackage[shortlabels]{enumitem}
\usepackage{tikz}
\newtheorem{theorem}{Theorem}
\DeclareMathOperator{\cas}{CAS}
\title{Distribution of independent sets in perfect $r$-ary trees}
\author{Daniel I\v{l}kovi\v{c}, Jun Yan}
\date{Source: arXiv:2601.16953v1 (CC BY 4.0)}
\begin{document}
\maketitle

\noindent\emph{Source and licence.} Distribution of independent sets in perfect $r$-ary trees. Version arXiv:2601.16953v1 (CC BY 4.0), distributed by arXiv under the Creative Commons Attribution 4.0 International licence, http://creativecommons.org/licenses/by/4.0/, as recorded in the arXiv licence metadata for this exact version and shown on the abstract page https://arxiv.org/abs/2601.16953v1. Full licence notice: \texttt{licenses/arxiv-2601.16953-CC-BY-4.0.txt}.\par\smallskip

\noindent\textit{Editorial scope.} Counted unit: the article's theorem thm:rarity (source0.tex lines 744--746). The article's complete original proof of it is delivered with it (source0.tex lines 748--768, one \texttt{proof} environment of 21 lines). No mathematics is written, repaired or re-derived: the statement and the proof are the article's own lines, in the article's own notation, and no retained line is substituted or re-flowed.  Retained uncounted as the source-local context the proof needs: lines 536--536; lines 615--615.  Nothing was omitted from any retained span, so the package carries no omission record.  No retained line contains a citation, so the wrapper carries no bibliography and no bibliography entry.  No retained line contains a figure environment, an image-inclusion command or drawing code, so no image is delivered and the package declares no asset dependency.  Editorial formatting only, in addition: an AMS \texttt{amsart} preamble that restates the article's own declarations and macros used by the retained lines, namely the environments \texttt{theorem} on separate counters, together with the standard packages those lines need.  The retained environments print in this document as Theorem 1 in order of appearance, whereas the article prints the same items with the numbering of its own full document.  The complete native source of the article as distributed in the arXiv e-print for this exact version is bundled unchanged as \texttt{sources/193281/source0.tex}.
\subsection{Case 1: $d(v,\ell)$ is even}\label{sec:even}

\subsection{Case 2: Odd Distance $d(v,\ell)$}\label{sec:odd}

\begin{theorem}\label{thm:rarity}
Let $T$ be a perfect $r$-ary tree with $r \geq 2$. For any integer $k \ge 1$, any vertex $v \in V(T)$, and any leaf $l \in V(T)$, we have $|\mathcal{I}_{T}^{k}(v)| \le |\mathcal{I}_{T}^{k}(l)|$.
\end{theorem}

\begin{proof}
The proof follows the strategy of the binary case. Due to the symmetry of perfect trees, the size of the star centreed at $v$ depends only on its depth. We assume without loss of generality that $l$ is a descendant of $c_1(v)$ and that the path $P_{v,\ell}$ consists entirely of edges connecting a parent to its first child. We construct an injective mapping $\phi_r: \mathcal{A}_{v,\ell} \to \mathcal{B}_{v,\ell}$ by generalizing the algorithms $\phi_{\textup{even}}$ and $\phi_{\textup{odd}}$.

We define the \textbf{Generalised Conditional Alternating Swap (Generalised CAS)}. Let $u$ and $d$ be two vertices on the path $P_{v,\ell}$. In the binary case, a shift along the path required swapping configurations between the single right subtrees $T_{R(u)}$ and $T_{R(d)}$. In the $r$-ary case, we must handle the sets of off-path siblings $S(u) = \{c_2(u), \dots, c_r(u)\}$ and $S(d) = \{c_2(d), \dots, c_r(d)\}$. Since $T$ is perfect, for each $j \in \{2, \dots, r\}$, the subtree rooted at $c_j(u)$ is isomorphic to the subtree rooted at $c_j(d)$. The Generalised CAS applies the binary CAS operation pairwise to these isomorphic subtrees. Specifically, if $I$ is the current independent set, the update is:
\[
\forall j \in \{2, \dots, r\}, \quad I \leftarrow \cas(I, T_{c_j(u)}, T_{c_j(d)}).
\]
Since the subtrees are disjoint, these operations are independent and preserve injectivity.

\textbf{Case 1: $d(v,\ell)$ is even.}
We define the mapping $\phi_{r, \textup{even}}$ analogously to Section~\ref{sec:even}. Let $I \in \mathcal{A}_{v,\ell}$.
\begin{enumerate}
    \item Initialise $I_1 = (I \setminus \{v\}) \cup \{\ell\}$. Maintain pointers $d_i$ (moving up from $p(\ell)$) and $u_i$ (moving down from $c_1(v)$).
    \item Perform the \textbf{Primary Check} and \textbf{Primary Shift} along $P_{v,\ell}$. If $d_i \in I_i$, shift the membership to $u_i$.
    \item Perform the \textbf{Generalised CAS} to swap the configurations of the sibling sets $S(d_i)$ and $S(u_i)$. This resolves the conflicts generated by adding $u_i$ and utilises the space freed by removing $d_i$.
    \item Perform the \textbf{Secondary Check} and \textbf{Secondary Shift} on the path vertices $c_1(u_i)$ and $p(d_i)$.
\end{enumerate}

\textbf{Case 2: $d(v,\ell)$ is odd.}
We define $\phi_{r, \textup{odd}}$ utilizing the synchronised propagation method from Section~\ref{sec:odd}. Initialise $Q_i^s$ with $\{p(\ell)\}$ and $Q_i^t$ with the set of all $r-1$ siblings of $c_1(v)$, i.e., $S(v) = \{c_2(v), \dots, c_r(v)\}$. The BFS propagation matches conflicts from the source to the target. When a node $x$ is processed (removed from the set), all its $r$ children $\{c_1(x), \dots, c_r(x)\}$ are added to the respective queue. Since $r \ge 2$, the capacity of the target queue is sufficient to resolve all potential conflicts, ensuring the mapping is well-defined and injective.
\end{proof}

\end{document}
